\section{Right-Generator Expansion Details}
\label{app:right-generator}

This appendix records the finite-order identity used by the certificate.  Let
\begin{equation}
 Z(t)=\log S_4(t)=tA+t^5B+t^7C+O(t^9).
\end{equation}
Using
\begin{equation}
 \operatorname{dexp}_{Z}(Z')=
 Z'+\frac12[Z,Z']+\frac16[Z,[Z,Z']]
 +\frac1{24}[Z,[Z,[Z,Z']]]+\cdots,
\end{equation}
and collecting powers through $t^7$ gives
\begin{align}
 S_4'(t)S_4(t)^{-1}-A
  ={}&5t^4B+2t^5[A,B]
  +t^6\left(7C+\frac23[A,[A,B]]\right)\nonumber\\
  &+t^7\left(3[A,C]+\frac16[A,[A,[A,B]]]\right)
  +R_{\geq8}(t).
 \label{eq:exact-right-expansion}
\end{align}
For example, the $t^5[A,B]$ coefficient comes from
$\frac12[tA,5t^4B]+\frac12[t^5B,A]=2t^5[A,B]$.
The later coefficients follow by the same homogeneous bookkeeping.  The
implementation constructs the series in the word algebra and regression-tests
\cref{eq:exact-right-expansion}; it does not substitute the printed identity
as an unchecked constant table.

The global contribution of a monomial $t^jD_j$ follows from
\cref{eq:global-duhamel}:
\begin{equation}
 r\int_0^{T/r}s^j\opnorm{D_j}\,ds
 =\frac{T^{j+1}\opnorm{D_j}}{(j+1)r^j}.
 \label{eq:monomial-global}
\end{equation}
This explains the degree labels in the result ledger and provides a direct
dimensional check on the inverse powers of $r$.

\section{Local Heisenberg Pauli Audit}
\label{app:local-pauli}

The single-bond lemma can be checked without symbolic software.  Let
$Q_1=XX$, $Q_2=YY$, and $Q_3=ZZ$.  For a phase-free two-qubit Pauli
$P=P_1\otimes P_2$, the number of $Q_j$ that anticommute with $P$ is either
zero or two.  One compact enumeration is shown in \cref{tab:local-audit}; rows
and columns give the one-qubit factors.

\begin{table}[ht]
\centering
\caption{Number of anticommuting terms among $XX,YY,ZZ$ for every local
two-qubit Pauli.}
\label{tab:local-audit}
\begin{tabular}{@{}c|rrrr@{}}
\toprule
$P_1\backslash P_2$ & $I$ & $X$ & $Y$ & $Z$ \\
\midrule
$I$ & 0 & 2 & 2 & 2 \\
$X$ & 2 & 0 & 2 & 2 \\
$Y$ & 2 & 2 & 0 & 2 \\
$Z$ & 2 & 2 & 2 & 0 \\
\bottomrule
\end{tabular}
\end{table}

If the count is zero, $[h,P]=0$.  If it is two, each contributing Hamiltonian
coefficient has magnitude $1/4$ and the Pauli commutator supplies magnitude
two, giving two output coefficients of magnitude $1/2$.  Their
Pauli-$\ell_1$ sum is one.  Hence
\begin{equation}
  \paulinorm{[h,P]}\leq1=\paulinorm{P}
\end{equation}
for a unit Pauli input, and linearity extends the stated growth majorant.  The
generic bound obtained by summing the three Hamiltonian Paulis independently
would give $3/2$; the exact parity enumeration removes the impossible third
contribution.

\section{Verifier Pseudocode}
\label{app:verifier}

\begin{algorithm}[ht]
\caption{Independent verification of the principal certificate}
\label{alg:verifier}
\begin{algorithmic}[1]
\Require main certificate $C$ and sidecar directory
\State validate schema version and required benchmark fields
\State recompute SHA-256 for every referenced sidecar
\State verify formula, normalization, $T$, $\eps$, and cost-model identities
\State load the D4 coefficient ledger and declared group partition
\State verify exact coverage: no missing, duplicate, or unknown Pauli key
\For{every group $G$}
  \For{every unordered pair $(P,Q)\subseteq G$}
    \State require symplectic parity$(P,Q)=1$
  \EndFor
  \State recompute the rational squared-weight sum
  \State require declared upper root squared $\geq$ recomputed sum
\EndFor
\State sum all group roots and compare to the D4 cell enclosure
\State verify D5--D7 and tail preconditions and contributions
\State sum the accepted-step global rational; require result $\leq\eps$
\State sum the previous-step rational; require result $>\eps$
\State recompute $G(r)$, bond propagators, CNOT totals, and exact ratio
\State return a structured success record; fail closed on any mismatch
\end{algorithmic}
\end{algorithm}

Deep mode inserts regeneration of the baseline center scan, complete formula
logarithm, DSW defects, Pauli coefficient map, and higher-degree majorants
before the certificate comparisons.  Importantly, both modes fail closed:
unknown schema versions, absent files, extra group members, and invalid
interval orientation are errors rather than warnings.

\section{Certificate Schema Sketch}
\label{app:schema}

\begin{longtable}{@{}p{0.29\textwidth}p{0.19\textwidth}p{0.44\textwidth}@{}}
\caption{Selected schema-v3 fields and verification obligations.}
\label{tab:schema}\\
\toprule
Field & Type & Obligation \\
\midrule
\endfirsthead
\toprule
Field & Type & Obligation \\
\midrule
\endhead
\path{benchmark} & object & Exact model, normalization, size, time,
tolerance, and primary metric match the claim \\
\path{published_baseline.site_density_upper} & integer pair & Denominator is
positive and deep regeneration is enclosed \\
\path{candidate.d4_certificate.path} & relative path & Resolves inside the
bundle; no path traversal \\
\path{candidate.d4_certificate.sha256} & 64-hex digest & Equals the bytes of
the resolved sidecar \\
\path{candidate.contributions} & rational pairs & Each pair has positive
denominator and agrees with recomputed bounds \\
\path{candidate.global_error_upper} & rational pair & Equals or dominates
the contribution sum and is below tolerance \\
\path{candidate.previous_step_error_upper} & rational pair & Correctly
evaluated at $r-1$ and exceeds tolerance \\
\path{cost_model} & integers/rule & Resource identities are evaluated without
floating point \\
\path{claimed_resources} & integer object & Every total is reconstructed from
the step count and cost model \\
\path{claims.exact_improvement_ratio} & integer pair & Equals the declared
baseline and candidate group counts \\
\bottomrule
\end{longtable}

\section{Additional Interpretation of the Improvement}
\label{app:interpretation}

Several ratios appear in the analysis and should not be conflated:
\begin{align}
\text{D4 norm tightening}
  &=\frac{20.160968407335066}{6.472926505087888}
    =3.114660484942633\ldots,\\
\text{step ratio}
  &=\frac{393}{97}=4.051546391752577\ldots,\\
\text{compiled resource ratio}
  &=\frac{30\cdot393+1}{30\cdot97+1}
    =\frac{11791}{2911}
    =4.050498110614909\ldots.
\end{align}
The step and resource ratios differ because the first and last fragment groups
merge to give the additive one in \cref{eq:merged-groups}.  The norm ratio is
not expected to equal either: a fourth-order global error scales roughly as
$r^{-4}$ only in a leading asymptotic regime, while the actual accepted point
contains D5--D7, a tail, and integer rounding.

The improvement relative to the challenge's twofold target is
\begin{equation}
 \frac{(11791/2911)}{2}=2.025249055307454\ldots.
\end{equation}
This means the achieved multiplier is just over twice the requested
multiplier; it does not mean an eightfold resource reduction.

\section{Checklist for Independent Review}
\label{app:checklist}

A reviewer can audit the result in increasing order of cost:
\begin{enumerate}
\item Read the benchmark tuple and confirm that candidate and control use the
same Hamiltonian, formula, tolerance, and cost model.
\item Run the manuscript headline validator and the package checksum check.
\item Run the fast verifier and inspect its structured output.
\item Confirm the exact integer map from 393/97 steps to all three resources.
\item Run the default unit suite, including mutation and local-Pauli tests.
\item Run D5 \code{--verify-only} to confirm the follow-up is separate and
marked noncertified at 78 steps.
\item If computational resources permit, run the deep verifier to regenerate
the formula-specific algebra and baseline center scan.
\item Rebuild this PDF and visually inspect figures, tables, citations, and
page boundaries; PDF appearance is a communication check, not a proof check.
\end{enumerate}
